% GENERATED FILE. Edit canonical lessons, then run npm run generate:books.
% canonical-source-hash: fnv1a64:1f8b8d311e5849ee
% canonical-lessons: BN-C495-beton, BN-C495-bonas, BN-C495-uparjon, BN-C495-penshon, BN-C495-poshano

\chapter{Pay Day}
\label{ch:bn-pay-day}
\hlchaptermodality{495}

\begin{chapteropening}
\textbf{By the end of this chapter:} \emph{I can talk about pay, bonuses and whether a job pays enough.}

Complete the last lesson of chapter 495: I can talk about pay, bonuses and whether a job pays enough.
\end{chapteropening}

\section[\texorpdfstring{\bn{বেতন} (betôn)}{বেতন (betôn)}]{\bn{বেতন} (betôn) — pay; wages}

\label{lesson:BN-C495-beton}



\begin{quote}
Before the new one: say the Bengali for a job interview, then the Bengali for to join a job.
\end{quote}

\subsection*{You'll want to know: \bn{বেতন}}
\textbf{\bn{বেতন}} — \emph{betôn} — \textquotedblleft{}pay; wages\textquotedblright{}.

\begin{quote}
\bn{মাসের} \bn{শেষে} \bn{আমরা} \bn{বেতন} \bn{পাই।}
\end{quote}

\begin{quote}
\emph{Māser sheshe āmrā betôn pāi.}
\end{quote}

\begin{quote}
We get our pay at the end of the month.
\end{quote}

\begin{cousinweb}
\textbf{\bn{বেতন}} (\emph{betôn}) is a tatsama, Sanskrit \emph{vetana}, \textquotedblleft{}wages, hire\textquotedblright{} (Monier-Williams).
\end{cousinweb}

\begin{grammarlens}[title={What you've built}]
One more word for describing what you lived and felt.
\end{grammarlens}

\subsection*{Your turn}
\begin{itemize}
\raggedright
  \item \emph{Say it:} \emph{betôn}
  \item \emph{Say it:} \emph{betôn}, once more
  \item \emph{Say it:} \emph{Māser sheshe āmrā betôn pāi.} — We get our pay at the end of the month.
  \item \emph{Recall:} say the Bengali for a job interview, then the Bengali for to join a job, then say \emph{betôn} again
\end{itemize}

\begin{tcolorbox}[breakable,colback=teal!4,colframe=teal!35!black,title={Before you move on}]
What does \bn{বেতন} mean? (\textquotedblleft{}Pay; wages\textquotedblright{}.) Say it once more.
\end{tcolorbox}
\section[\texorpdfstring{\bn{বোনাস} (bonās)}{বোনাস (bonās)}]{\bn{বোনাস} (bonās) — a bonus}

\label{lesson:BN-C495-bonas}



\begin{quote}
Before the new one: say the Bengali for to join a job, then the Bengali for pay.
\end{quote}

\subsection*{You'll want to know: \bn{বোনাস}}
\textbf{\bn{বোনাস}} — \emph{bonās} — \textquotedblleft{}a bonus\textquotedblright{}.

\begin{quote}
\bn{এই} \bn{বছর} \bn{আমরা} \bn{বোনাস} \bn{পেলাম।}
\end{quote}

\begin{quote}
\emph{Ei bôchhôr āmrā bonās pelām.}
\end{quote}

\begin{quote}
This year we got a bonus.
\end{quote}

\begin{cousinweb}
\textbf{\bn{বোনাস}} (\emph{bonās}) is borrowed from English \emph{bonus}.
\end{cousinweb}

\begin{grammarlens}[title={What you've built}]
One more word for describing what you lived and felt.
\end{grammarlens}

\subsection*{Your turn}
\begin{itemize}
\raggedright
  \item \emph{Say it:} \emph{bonās}
  \item \emph{Say it:} \emph{bonās}, once more
  \item \emph{Say it:} \emph{Ei bôchhôr āmrā bonās pelām.} — This year we got a bonus.
  \item \emph{Recall:} say the Bengali for to join a job, then the Bengali for pay, then say \emph{bonās} again
\end{itemize}

\begin{tcolorbox}[breakable,colback=teal!4,colframe=teal!35!black,title={Before you move on}]
What does \bn{বোনাস} mean? (\textquotedblleft{}A bonus\textquotedblright{}.) Say it once more.
\end{tcolorbox}
\section[\texorpdfstring{\bn{উপার্জন} (upārjôn)}{উপার্জন (upārjôn)}]{\bn{উপার্জন} (upārjôn) — earnings; what one earns}

\label{lesson:BN-C495-uparjon}



\begin{quote}
Before the new one: say the Bengali for pay, then the Bengali for a bonus.
\end{quote}

\subsection*{You'll want to know: \bn{উপার্জন}}
\textbf{\bn{উপার্জন}} — \emph{upārjôn} — \textquotedblleft{}earnings; what one earns\textquotedblright{}.

\begin{quote}
\bn{দোকান} \bn{থেকে} \bn{তার} \bn{উপার্জন} \bn{ভালো।}
\end{quote}

\begin{quote}
\emph{Dokān theke tār upārjôn bhālo.}
\end{quote}

\begin{quote}
Her earnings from the shop are good.
\end{quote}

\begin{cousinweb}
\textbf{\bn{উপার্জন}} (\emph{upārjôn}) is a tatsama, Sanskrit \emph{upārjana}, \textquotedblleft{}the act of gaining\textquotedblright{} (Monier-Williams).
\end{cousinweb}

\begin{grammarlens}[title={What you've built}]
One more word for describing what you lived and felt.
\end{grammarlens}

\subsection*{Your turn}
\begin{itemize}
\raggedright
  \item \emph{Say it:} \emph{upārjôn}
  \item \emph{Say it:} \emph{upārjôn}, once more
  \item \emph{Say it:} \emph{Dokān theke tār upārjôn bhālo.} — Her earnings from the shop are good.
  \item \emph{Recall:} say the Bengali for pay, then the Bengali for a bonus, then say \emph{upārjôn} again
\end{itemize}

\begin{tcolorbox}[breakable,colback=teal!4,colframe=teal!35!black,title={Before you move on}]
What does \bn{উপার্জন} mean? (\textquotedblleft{}Earnings; what one earns\textquotedblright{}.) Say it once more.
\end{tcolorbox}
\section[\texorpdfstring{\bn{পেনশন} (penshôn)}{পেনশন (penshôn)}]{\bn{পেনশন} (penshôn) — a pension}

\label{lesson:BN-C495-penshon}



\begin{quote}
Before the new one: say the Bengali for a bonus, then the Bengali for earnings.
\end{quote}

\subsection*{You'll want to know: \bn{পেনশন}}
\textbf{\bn{পেনশন}} — \emph{penshôn} — \textquotedblleft{}a pension\textquotedblright{}.

\begin{quote}
\bn{দাদু} \bn{প্রতি} \bn{মাসে} \bn{পেনশন} \bn{পান।}
\end{quote}

\begin{quote}
\emph{Dādu prôti māse penshôn pān.}
\end{quote}

\begin{quote}
My grandfather gets a pension every month.
\end{quote}

\begin{cousinweb}
\textbf{\bn{পেনশন}} (\emph{penshôn}) is borrowed from English \emph{pension}.
\end{cousinweb}

\begin{grammarlens}[title={What you've built}]
One more word for describing what you lived and felt.
\end{grammarlens}

\subsection*{Your turn}
\begin{itemize}
\raggedright
  \item \emph{Say it:} \emph{penshôn}
  \item \emph{Say it:} \emph{penshôn}, once more
  \item \emph{Say it:} \emph{Dādu prôti māse penshôn pān.} — My grandfather gets a pension every month.
  \item \emph{Recall:} say the Bengali for a bonus, then the Bengali for earnings, then say \emph{penshôn} again
\end{itemize}

\begin{tcolorbox}[breakable,colback=teal!4,colframe=teal!35!black,title={Before you move on}]
What does \bn{পেনশন} mean? (\textquotedblleft{}A pension\textquotedblright{}.) Say it once more.
\end{tcolorbox}
\section[\texorpdfstring{\bn{পোষানো} (poshāno)}{পোষানো (poshāno)}]{\bn{পোষানো} (poshāno) — to be worth it; to pay enough}

\label{lesson:BN-C495-poshano}



\begin{quote}
Before the new one: say the Bengali for earnings, then the Bengali for a pension.
\end{quote}

\subsection*{You'll want to know: \bn{পোষানো}}
\textbf{\bn{পোষানো}} — \emph{poshāno} — \textquotedblleft{}to be worth it; to pay enough\textquotedblright{}.

\begin{quote}
\bn{এই} \bn{মাইনেতে} \bn{আমার} \bn{পোষায়} \bn{না।}
\end{quote}

\begin{quote}
\emph{Ei māinete āmār poshāy nā.}
\end{quote}

\begin{quote}
This salary isn't worth my while.
\end{quote}

Past (he/she …): \bn{সেই} \bn{মাইনেতে} \bn{তার} \bn{পোষাল} \bn{না}, \emph{sei māinete tār poshālô nā}.

The person it pays takes the genitive: \bn{আমার} \bn{পোষায়} \bn{না}, it isn't worth it for me.

\begin{cousinweb}
\textbf{\bn{পোষানো}} (\emph{poshāno}) is a tadbhava from Sanskrit \emph{poṣayati}, \textquotedblleft{}nourishes, keeps\textquotedblright{} (CDIAL 8410); the sense \textquotedblleft{}to be worth it\textquotedblright{} is Bengali's own.
\end{cousinweb}

\begin{grammarlens}[title={What you've built}]
One more word for describing what you lived and felt.
\end{grammarlens}

\subsection*{Your turn}
\begin{itemize}
\raggedright
  \item \emph{Say it:} \emph{poshāno}
  \item \emph{Say it:} \emph{poshāno}, once more
  \item \emph{Say it:} \emph{Ei māinete āmār poshāy nā.} — This salary isn't worth my while.
  \item \emph{Recall:} say the Bengali for earnings, then the Bengali for a pension, then say \emph{poshāno} again
\end{itemize}

\begin{tcolorbox}[breakable,colback=teal!4,colframe=teal!35!black,title={Before you move on}]
What does \bn{পোষানো} mean? (\textquotedblleft{}To be worth it; to pay enough\textquotedblright{}.) Say it once more.
\end{tcolorbox}
